% !TEX root = ../Thesis.tex
\section{The Trial, and the Question That Survives Its Conclusion}

MASTER DAPT randomized 4,579 high bleeding risk patients, after coronary
stent implantation, between abbreviated and prolonged dual antiplatelet
therapy. Randomization took place one month after the index procedure, in
patients who were free of events up to that point: \emph{that visit is the
temporal landmark of the whole analysis}, and every covariate used in this
thesis is measured at or before it.

The trial concluded in favour of the abbreviated strategy — noninferior on
ischemic events, superior on bleeding. That is an average result, and it is
the concrete instance of the problem posed in Chapter 1: a favourable ATE is
perfectly compatible with subgroups that gain nothing, or that are actively
harmed, and the trial on its own does not say which of the two this cohort
contains. The question that survives the trial's conclusion is therefore not
\emph{whether} abbreviating is better on average — that is settled — but
\emph{for whom} it is better, and by how much.

\section{The Cohort}

The analysis dataset contains the 4,579 randomized patients, split almost
exactly in half by the randomization (2,295 abbreviated, 2,284 prolonged),
described by 63 baseline covariates plus the treatment indicator. The
covariates were selected from the full trial case report form through a
combination of predictive screening and consultation with the study domain
experts, and cover demographic, clinical, laboratory and procedural
information: 15 are continuous measurements (age, laboratory values such as
creatinine, haemoglobin and troponin, procedural quantities such as contrast
volume) and the remaining 48 are binary or categorical indicators of
comorbidities, concomitant medications and procedural characteristics.

Missing values are deliberately \emph{not} imputed in the shared dataset.
Eight covariates carry missing entries — most notably pre-procedural glycemia,
platelet count and creatinine — and every imputation is deferred to the
individual modelling pipelines, where it is fitted inside the
cross-validation folds. Imputing once, up front, on the full dataset would
leak information across folds and inflate every performance estimate
downstream; deferring it costs some convenience and buys the guarantee that
the out-of-fold numbers reported in Chapters 5 to 7 mean what they claim to.

\paragraph{Sign convention.}
The coding of the treatment is fixed for the entire thesis: $T = 1$ denotes
\emph{prolonged} therapy and $T = 0$ denotes \emph{abbreviated} therapy. With
outcomes coded as adverse events, the conditional average treatment effect
\[
  \tau(x) \;=\; \mathbb{E}\!\left[\,Y(1) - Y(0) \mid X = x\,\right]
\]
is therefore the excess event risk caused by prolonging, which is exactly the
\emph{benefit of abbreviating}. A positive $\tau(x)$ means that patients with
profile $x$ are better off with the abbreviated strategy; a negative $\tau(x)$
means they are better off with the prolonged one; a value near zero means the
choice is immaterial for that patient type. This convention is used without
exception in every chapter that follows, so that the sign of an estimate can
always be read directly as \emph{abbreviate} or \emph{do not abbreviate}.

\paragraph{Endpoints.}
Of the nine adjudicated binary endpoints available at 335 days, five are
carried through the thesis; they are reported in Table~\ref{tab:endpoints}.
Time-to-event information is deliberately discarded in favour of the binary
status at 335 days, which is the horizon of the trial's own primary analysis:
follow-up is administratively complete at that point and essentially uniform
across patients, so the survival machinery would add complexity without
adding information. These five endpoints are the targets of the risk models
of Chapter 5 and of the effect estimates of Chapters 6 and 7.

\begin{table}[h]
\centering
\caption{Binary endpoints at 335 days.}
\label{tab:endpoints}
\begin{tabular}{lrr}
\toprule
Endpoint & Events & Rate \\
\midrule
Bleeding                    & 506 & 11.1\% \\
Bleeding (BARC 2/3/5)       & 359 &  7.8\% \\
Myocardial infarction (MI)  & 109 &  2.4\% \\
Cardiovascular death        &  81 &  1.8\% \\
Stroke                      &  35 &  0.8\% \\
\bottomrule
\end{tabular}
\end{table}

The two bleeding endpoints are by far the most frequent, at 11.1\% and 7.8\%
respectively; the BARC 2/3/5 definition is the more restrictive of the two,
which is why it records fewer events than overall bleeding — the difference
consists of the milder BARC 1 episodes, which do not require medical
attention. The three ischemic endpoints follow well below, with myocardial
infarction at 2.4\%, cardiovascular death at 1.8\% and stroke at 0.8\%.

Table~\ref{tab:endpoints} also makes explicit how imbalanced these endpoints
are: event counts span more than an order of magnitude, from 506 bleeding
events down to 35 strokes, and three of the five endpoints — myocardial
infarction, cardiovascular death and stroke — fall below 110 events overall.
This imbalance is not incidental: because MASTER DAPT enrolled exclusively
high bleeding risk patients, statistical power is concentrated on the bleeding
side of the trade-off, while the rarer ischemic endpoints leave comparatively
little signal to work with. This limitation resurfaces in later chapters, both
in how reliably each endpoint can be predicted (Chapter 5) and in how
confidently treatment effect heterogeneity can be detected on it
(Chapters 6--7).

\paragraph{Two facts stated here because they return later.}
First, \emph{every} patient in this cohort is at high bleeding risk, because
that was an eligibility criterion and not an accident of the randomization.
The cohort is therefore a deliberately restricted slice of the PCI population,
not a sample of it: the spread of bleeding risk within it is compressed by
construction. Chapter 7 shows that this range restriction, and not a
modelling failure, is the reason a bleeding risk model that is perfectly
sensible in itself discriminates far less well here than the published
literature would suggest.

Second, ischemic events are an order of magnitude rarer than bleeding events,
with cardiovascular death and stroke both below one hundred observations. Any
heterogeneity in the ischemic arm of the trade-off has to be detected from
that handful of events, which caps in advance how much can honestly be claimed
about it — a constraint that Chapter 7 makes quantitative.

\section{The \texttt{fup\_} Variables Are Not Leakage}

Three of the 63 covariates carry the prefix \texttt{fup\_} —
\texttt{fup\_other}, \texttt{fup\_nitrates} and \texttt{fup\_insulin} — which
suggests \emph{follow-up}, and therefore suggests information recorded after
treatment assignment. Used as predictors, such variables would be leakage:
they would let a model see part of what it is supposed to predict, and would
make every performance and effect estimate optimistic in a way no
cross-validation scheme can repair.

They are not leakage here. Randomization in MASTER DAPT occurs at the
one-month visit, and that visit \emph{is} the landmark: the \texttt{fup\_}
prefix refers to the follow-up schedule of the trial's case report form, but
the values in question are recorded at that visit, before treatment assignment
takes effect. They are consequently baseline covariates in the sense that
matters — measured before the intervention whose effect is being estimated.

The point is worth the space it occupies for two reasons. It fixes, once, the
criterion applied to all 63 covariates: what determines admissibility is the
recording date relative to the landmark, not the prefix of the variable name,
and a mechanical rule based on names would in this dataset discard valid
information while retaining nothing extra. And it closes in advance the most
obvious objection to the negative result of Chapters 6 and 7 — namely that the
absence of detectable heterogeneity is an artefact of having thrown away good
predictors.

\section{A First Look at the Data}

Before any modelling, two exploratory questions are worth settling: whether
the cohort visibly separates into subgroups, and whether obvious pairwise
structure among the covariates exists.

\paragraph{Structure and clustering.}
Principal component analysis on the standardized covariates shows no dominant
axis: a substantial number of components is needed to reach 80--90\% of the
variance, indicating that the covariates are only weakly correlated with each
other. The leading component of the continuous block is dominated by kidney
function markers (pre- and post-procedural creatinine) and myocardial injury
markers (troponin), reflecting the expected comorbidity pattern of a PCI
population. Projecting patients onto the leading components and colouring them
by outcome reveals no sharp clustering for any endpoint, and the nonlinear
embeddings agree: both t-SNE and UMAP, applied to the leading principal
components, scatter event patients across the entire manifold with no distinct
region for any endpoint. Treatment assignment likewise shows no separation,
which is the expected signature of a balanced randomization. Bleeding
endpoints are marginally better separated than ischemic ones, which is the
same asymmetry that Table~\ref{tab:endpoints} predicts and that Chapter 5
recovers quantitatively.

\includegraphics[width=1\textwidth]{Non_linear_embeddings.png}

The absence of clustering also rules out a simpler story for the ischemic
endpoints: that events concentrate in a small set of multivariate extremes.
Flagging patients by Hotelling's $T^2$ statistic on the leading ten
components identifies 162 of 4,579 patients (3.5\%) as extreme outliers at
the 99.5\% level, and these outliers do carry a higher event rate — death and
stroke occur roughly 3.4 and 3.5 times more often among them than among the
rest of the cohort. But the outlier set is too small for that elevated rate
to matter in aggregate: only 11.4\% of cardiovascular deaths and 11.6\% of
strokes fall among the outliers, meaning that close to nine in ten ischemic
events occur in patients whose covariate profile is otherwise unremarkable.
Whatever drives ischemic risk in this cohort, it is not concentrated in a
handful of biomarker extremes.

\includegraphics[width=1\textwidth]{Outliers_analysis.png}

\paragraph{Pairwise interactions.}
A penalized search over pairwise products of a curated set of clinically
interpretable covariates flags a few interaction terms that survive $L_1$
shrinkage and improve out-of-fold discrimination for bleeding. These are
recorded but not pursued, because they answer a different question from the
one this thesis asks: they are associations \emph{among covariates} in
predicting an outcome, not modifications of the \emph{treatment effect}. An
interaction between two baseline characteristics tells us the risk surface is
not additive; it says nothing about whether the benefit of abbreviating varies
across patients. That question requires the treatment to enter the
interaction, and it is the subject of Chapter 6, not of this one.

\includegraphics[width=1\textwidth]{Lasso_Regularization_Path.png}